% from ega2-1-fr.tex
\oldpage[II]{47}

\begin{env}[2.8.14]
\label{II.2.8.14-fr}
Conservant toujours les hypothèses et notations de
(\hyperref[II.2.8.9-fr]{2.8.9}), soit $N'$ un second $S'$-module gradué, et
soit $N=N'\otimes_{A'}A$. Il est immédiat que les di-homomorphismes
canoniques $M'\to M$, $N'\to N$ donnent un di-homomorphisme
$M'\otimes_{S'}N'\to M\otimes_S N$ (relatif à l'homomorphisme canonique
d'anneaux $S'\to S$), et par suite aussi un $S$-homomorphisme de degré~$0$
\[
  (M'\otimes_{S'}N')\otimes_{A'}A\longrightarrow M\otimes_S N,
\]
auquel correspond (compte tenu de
(\hyperref[II.2.8.10-fr]{2.8.10})) un homomorphisme de
$\mathcal{O}_X$-Modules
\[
\label{II.2.8.14.1-fr}
  \Phi^*(\widetilde{M'\otimes_{S'}N'})
  \longrightarrow\widetilde{M\otimes_S N}.
\tag{2.8.14.1}
\]

En outre, on vérifie aussitôt que le diagramme
\[
\label{II.2.8.14.2-fr}
  \xymatrix{
    \Phi^*(\widetilde{M'}\otimes_{\mathcal{O}_{X'}}\widetilde{N'})
      \ar[r]^{\sim}\ar[d]_{\Phi^*(\lambda)}
    & \widetilde{M}\otimes_{\mathcal{O}_X}\widetilde{N}
      \ar[d]^{\lambda}
    & =\Phi^*(\widetilde{M'})\otimes_{\mathcal{O}_X}
      \Phi^*(\widetilde{N'})\\
    \Phi^*(\widetilde{M'\otimes_{S'}N'})\ar[r]
    & \widetilde{M\otimes_S N}
  }
\tag{2.8.14.2}
\]
est commutatif, la première ligne étant l'isomorphisme canonique
(\hyperref[0.4.3.3-fr]{0, 4.3.3}). Si l'idéal $S'_+$ est engendré par
$S'_1$, il est clair que $S_+$ est engendré par $S_1$, et les deux flèches
verticales de (\hyperref[II.2.8.14.2-fr]{2.8.14.2}) sont alors des
isomorphismes (\hyperref[II.2.5.13-fr]{2.5.13})~; il en est donc de même de
(\hyperref[II.2.8.14.1-fr]{2.8.14.1}).

On a de même un di-homomorphisme canonique
$\operatorname{Hom}_{S'}(M',N')\to\operatorname{Hom}_S(M,N)$ en faisant
correspondre à un homomorphisme $u'$ de degré $k$ l'homomorphisme
$u'\otimes1$, qui est aussi de degré $k$~; on en déduit encore un
$S$-homomorphisme de degré~$0$
\[
  (\operatorname{Hom}_{S'}(M',N'))\otimes_{A'}A
  \longrightarrow\operatorname{Hom}_S(M,N)
\]
d'où un homomorphisme de $\mathcal{O}_X$-Modules
\[
\label{II.2.8.14.3-fr}
  \Phi^*(\widetilde{\operatorname{Hom}_{S'}(M',N')})
  \longrightarrow\widetilde{\operatorname{Hom}_S(M,N)}.
\tag{2.8.14.3}
\]

En outre, le diagramme
\[
  \xymatrix{
    \Phi^*(\widetilde{\operatorname{Hom}_{S'}(M',N')})\ar[r]
      \ar[d]_{\Phi^*(\mu)}
    & \widetilde{\operatorname{Hom}_S(M,N)}\ar[d]^{\mu}\\
    \Phi^*(\mathcal{H}om_{\mathcal{O}_{X'}}
      (\widetilde{M'},\widetilde{N'}))\ar[r]
    & \mathcal{H}om_{\mathcal{O}_X}(\widetilde{M},\widetilde{N})
  }
\]
est commutatif (la seconde ligne horizontale étant l'homomorphisme canonique
(\hyperref[0.4.4.6-fr]{0, 4.4.6})).
\end{env}

